% frmcs_appsession_en.tex  — Draft v2.9
% XeLaTeX — ETSI FRMCS Application Session Concept (English)
\documentclass[11pt,a4paper]{article}

\usepackage{fontspec}
\setmainfont[SmallCapsFont = Avenir LT Std 35 Light]{Avenir LT Std 35 Light}
\setsansfont{Avenir LT Std 45 Book}

\usepackage{geometry}\geometry{a4paper,margin=2.5cm}
\usepackage{polyglossia}\setdefaultlanguage{english}
\usepackage{microtype,booktabs,longtable,array,listings,hyperref,titlesec,enumitem,soul,parskip}
\usepackage[table]{xcolor}
\usepackage{pifont}
\newcommand{\cmark}{\ding{51}}

\definecolor{codebg}{RGB}{245,245,245}
\definecolor{etsiblue}{RGB}{0,70,127}
\definecolor{noteblue}{RGB}{230,240,255}

\lstset{backgroundcolor=\color{codebg},basicstyle=\small\ttfamily,breaklines=true,
  frame=single,framerule=0.4pt,rulecolor=\color{gray!40},
  xleftmargin=6pt,xrightmargin=6pt,aboveskip=8pt,belowskip=4pt,
  showstringspaces=false,keepspaces=true}

\titleformat{\section}{\large\bfseries\sffamily\color{etsiblue}}{{\thesection}}{1em}{}
\titleformat{\subsection}{\normalsize\bfseries\sffamily\color{etsiblue}}{{\thesubsection}}{1em}{}
\titleformat{\subsubsection}{\normalsize\bfseries\sffamily}{{\thesubsubsection}}{1em}{}

\newenvironment{notebox}{%
  \begin{center}\begin{tabular}{|p{0.93\textwidth}|}\hline\rowcolor{noteblue}
}{\\\hline\end{tabular}\end{center}}

% Interface name shorthand
\newcommand{\OBAPP}{OB\textsubscript{APP}}
\newcommand{\TSAPP}{TS\textsubscript{APP}}
% Arrow substitute: robust in text and table contexts
\newcommand{\arr}{$\to$}

\newcommand{\svcref}{\texttt{svcSessionRef}}
\newcommand{\appcall}{\texttt{appSessionId\_call}}
\newcommand{\appterm}{\texttt{appSessionId\_term}}

\hypersetup{colorlinks=true,linkcolor=etsiblue,urlcolor=etsiblue,
  pdfauthor={R2DATO D25.3},pdftitle={ETSI FRMCS -- Application Session Concept}}

\begin{document}

\begin{titlepage}
  \vspace*{3cm}{\sffamily\color{etsiblue}
    {\Huge\bfseries ETSI FRMCS\\[0.4em]Application Session Concept}\\[1.5em]
    {\Large Proposed Modifications to ETSI TS 103\,764, TS 103\,765-3/4/5}\\[2em]
    {\large Working Contribution -- Draft v2.9}\\[0.5em]
    {\normalsize R2DATO D25.3 -- Concept Paper FRMCS}\\[4em]
    \hrule\vspace{0.5em}
    {\small\textbf{Target documents:} ETSI TS 103\,764 (RT-0095), TS 103\,765-3 (RT-0097),
      TS 103\,765-4 (RT-0099), TS 103\,765-5 (RT-00104)\\[0.3em]
      \textbf{Unchanged documents:} All UIC specifications (FFFIS-7950, FIS-7970,
      FRS FU-7120, AT-7800, TOBA FRS-7510, EIRENE FRS, EIRENE SRS)}}
\end{titlepage}

\tableofcontents\clearpage

% ── 1. ARCHITECTURAL PRINCIPLE ──────────────────────────────────────────────
\section{Architectural Principle}

\subsection{Reformulation of the Session Concept}

In the ETSI TSs, \textbf{"session"} previously referred to the MCData IPcon tunnel
established by the Gateway on behalf of an application (what this proposal calls a
\textbf{service session}). This proposal redefines the terms:

\begin{itemize}[leftmargin=2em]
  \item \textbf{"session"}: the application session -- communication between two FRMCS
    railway applications, composed of one or more service sessions managed internally
    by the Gateways.
  \item \textbf{"service session"}: an individual communication bearer. The first is
    necessarily MCx (MCData IPcon) -- unchanged. Additional ones may be any MCx type
    or future railway critical services (R2DATO D25.3 Part B); not limited to the same
    service domain.
\end{itemize}

All internal MCx tunnel occurrences of ``session'' in the ETSI TSs are replaced by
``service session''. Application-visible identifiers (\texttt{sessionId},
\texttt{POST /sessions}, \texttt{openSession*}, \texttt{closeSession*},
\texttt{incomingSessionNotif*}, \texttt{open session}) are \textbf{unchanged}.

\begin{notebox}
  \textbf{Nothing changes for applications.} \OBAPP\ and \TSAPP\ interfaces are
  strictly unchanged. UIC specs already use ``session'' in the application sense --
  their terminology is now naturally aligned without modification.
\end{notebox}

\subsection{Why ASAS Rather than SIP INVITE AnyExt}

RT-0097 §6.2.2.3 already carries \texttt{<application-data>} in the \texttt{AnyExt}
field (format ``not specified in the present document''). While usable for MCx, this
approach is structurally excluded: a future non-MCx critical service has no SIP INVITE
to carry the negotiation. The protocol-agnostic ASAS is the only mechanism that
survives the application/service strata separation of R2DATO D25.3 Part B.

\subsection{Application Identifier Resolution (AIR)}

The AIR is an internal Gateway function resolving \texttt{remoteId} to a
\textbf{service identifier} (MCx or future critical service).
\textbf{First-edition rules -- fully specified:}

\begin{enumerate}[leftmargin=2em]
  \item If \texttt{remoteId} ends with \texttt{.ertms} (ETCS format
    \texttt{id.ty.cc.ertms} as specified in UIC AT-7800 §5.3, which references
    UNISIG SUBSET-037): append the current MCx domain
    \arr\ \texttt{id.ty.cc.ertms@current\_domain}.
    Consistent with current UIC specifications; functional baseline (note: non-optimal
    in domain-overlap zones).
  \item Otherwise: pass-through as MCx FA -- identical to historical behaviour.
\end{enumerate}

\begin{notebox}
  The AIR function is designed as an \textbf{extension point}. Future editions may
  introduce DNS-based or FROP-configured resolution without any change to \OBAPP,
  \TSAPP, SIP, or UIC specifications, in connection with ERJU work on ETCS addressing
  in FRMCS (UIC EIRENE~SRS, UNISIG SUBSET-037).
\end{notebox}

\subsection{Initialisation Mechanism}

Once the first service session is established (SIP + GRE unchanged), each Gateway
exposes an \textbf{ASAS} reachable on its post-NAT IP address (SD\_PUB\_IP,
cf.\ §\ref{sec:nat}). The calling Gateway drives:

\begin{enumerate}[leftmargin=2em]
  \item \textbf{Initialisation}: POST to ASAS with \texttt{svcSessionId\_1} as
    correlator; receives \appterm, Ed25519 public key, alias list, MSP policy,
    keepalive period.
  \item \textbf{Additional service sessions}: established per alias list + FROP; each
    declared to ASAS with signed \texttt{sessionToken}, enabling mutual authentication
    independently of the underlying signalling protocol.
  \item \textbf{Keepalive}: sent only if applicative silence $\geq$ T\_KA on a given
    \svcref\ (intelligent keepalive, cf.\ §\ref{sec:ka}).
  \item \textbf{Alias list update}: PATCH sent immediately by terminating GW on any
    new alias or service identifier registration. No configuration required.
  \item \textbf{Release}: DELETE to ASAS + release each service session per its
    signalling protocol.
\end{enumerate}

\subsection{Application Session Lifecycle}

A session is alive while at least one service session is \texttt{ACTIVE} or
\texttt{DEGRADED}. When all are lost, a \textbf{recovery timer} starts:
INTEGER(1..180)\,s (FROP); SHALL be $\leq 5$\,s for sessions carrying ETCS/ATP
(AT-7800 Annex A.6). On expiry, the \textbf{application session} is released.

% ── 2. KEY CONCEPTS ─────────────────────────────────────────────────────────
\section{Key Concepts}

\subsection{Terminology}

\begin{longtable}{>{\bfseries}p{3.8cm} p{10.5cm}}
  \toprule Term & Definition \\ \midrule \endhead
  session & Application-level communication between two FRMCS railway applications,
    composed of one or more service sessions. Identified externally by \texttt{sessionId}
    (\OBAPP/\TSAPP), internally by (\appcall, \appterm). \\[0.3em]
  service session & Individual communication bearer, transparent to applications.
    First is necessarily MCx (MCData IPcon). Additional may be MCx or future critical
    services. Identified by \svcref\ within its parent session. \\[0.3em]
  open session & A session that has not been released (unchanged). \\[0.3em]
  AIR & Internal Gateway function resolving \texttt{remoteId} to a service identifier
    (MCx or future critical service). First-edition rules (fully specified):
    \texttt{*.ertms} \arr\ \texttt{*.ertms@current\_domain} (UIC AT-7800~§5.3,
    ref.\ UNISIG SUBSET-037); otherwise pass-through FA MCx.
    Extension point for future optimisation (DNS/FROP, in connection with ERJU work
    on ETCS addressing -- UIC EIRENE~SRS, UNISIG SUBSET-037).
    Not visible to applications. \\[0.3em]
  ASM & Internal Gateway function managing session and service session lifecycle,
    driving ASAS interactions and MSP policy. Not visible to applications. \\[0.3em]
  ASAS & REST API server exposed by each Gateway on its post-NAT IP (SD\_PUB\_IP).
    Security is based on per-message Ed25519 authentication; there is no normative
    requirement for confidentiality of ASAS exchanges between Gateways. HTTP transport
    MAY be protected by TLS (recommended). Ed25519 key exchange anchored on first MCx
    service session -- no inter-GW PKI required. \\[0.3em]
  MSP & Multi Service Path: inter-Gateway mechanism managing multiple parallel service
    sessions within a single application session, operated by the ASMs. Distinct from
    transport-level Multipath (RT-00104). Policy (\texttt{mspPolicy}) negotiated at
    ASAS initialisation: DUPLICATION, RECOVERY or LOADBALANCE. \\[0.3em]
  Calling GW & GW initiating \texttt{POST /sessions} via \OBAPP.
    Derives ASAS IP from post-NAT SIP source IP. \\[0.3em]
  Terminating GW & GW hosting destination application. Exposes ASAS. Builds alias
    list; SHALL immediately PATCH on any new alias or service identifier registration;
    notifies application via \TSAPP\ (\texttt{incomingSessionNotif*}). \\[0.3em]
  alias list & Dynamic list of service identifiers + types + domains.
    Returned at initialisation; immediately updated via \texttt{PATCH /aliases}
    on any new registration -- no FROP configuration required. \\[0.3em]
  service identifier & Identifier used to establish a service session, regardless
    of service type. For MCx: FA or MC Service User ID. For future critical services:
    any identifier the terminating Gateway can associate with the destination user. \\[0.3em]
  recovery timer & INTEGER(1..180)\,s (FROP); SHALL be $\leq 5$\,s for ETCS/ATP
    (AT-7800 Annex A.6). Application session released on expiry. \\[0.3em]
  T\_KA & Keepalive period: INTEGER(1..10)\,s (application profile).
    Rule: T\_KA $\leq$ recovery\_timer\,/\,2 ($\leq 2.5$\,s for ETCS/ATP).
    $N = 2$ consecutive failures (fixed) \arr\ \texttt{DEGRADED}. \\
  \bottomrule
\end{longtable}

\subsection{Session Identifiers}

\begin{itemize}[leftmargin=2em]
  \item \appcall\ (UUID v4): calling GW, stable for session lifetime.
  \item \appterm\ (UUID v4): terminating GW at ASAS init. Analogous to GTP TEIDs.
  \item \svcref\ (UUID v4): identifies a service session. \textbf{Stable and immutable},
    including during recovery.
  \item \texttt{svcSessionId} (native): SIP Call-ID for MCx. Used as correlator.
    May change during recovery.
\end{itemize}

\subsection{Authentication Model}
\label{sec:trust}

\textbf{The first MCx service session is a trusted service already responsible for
identifying the parties} -- its establishment is a sufficient authentication anchor.
There is no normative requirement for confidentiality of ASAS exchanges between
Gateways: the ASAS operates in a railway inter-GW environment where protection of
session metadata is not a normative prerequisite.

\subsubsection*{Level 1 -- HTTP Transport (TLS recommended)}

The ASAS HTTP transport MAY be protected by TLS~1.3 (recommended for transport
integrity). Server certificate may be \textbf{self-signed} -- no inter-operator PKI
required. TLS is not normative at this level; normative authentication is provided
by Levels~2 and~3. The mTLS option is possible if a bilateral OOB FSNNI agreement
so provides (FROP).

\begin{notebox}
  \textbf{FSNNI topology hiding note}: a potential intermediate domain performing
  topology hiding does not compromise authentication -- the Ed25519
  \texttt{sessionToken} provides per-message integrity regardless of the network
  path. The only constraint is ASAS reachability on SD\_PUB\_IP
  (covered by §\ref{sec:nat}).
\end{notebox}

\subsubsection*{Level 2 -- Authentication Anchor (First MCx Service Session)}

The MC Service Server authenticated the calling Gateway before accepting the session.
Ed25519 public key exchange over this channel bootstraps authentication: any key
received at ASAS initialisation is authentic by construction.

\subsubsection*{Level 3 -- Per-Message Authentication (Ed25519 sessionToken)}

\begin{lstlisting}[language={}]
sessionToken = Sign(sk_call,
  appSessionId_call || 0x00 || appSessionId_term || 0x00 ||
  svcSessionRef     || 0x00 || svcSessRequestTime || 0x00 ||
  svcSessAcceptTime)
// Recovery: append 0x00 || replSvcSessionRef
ackToken = Sign(sk_term, same input)
\end{lstlisting}

Serialisation: UTF-8 concatenation with \texttt{0x00} separators -- deterministic,
no external dependency.

\subsection{Intelligent Keepalive}
\label{sec:ka}

\begin{itemize}[leftmargin=2em]
  \item \textbf{T\_KA}: INTEGER(1..10)\,s (application profile);
    SHALL be $\leq 2.5$\,s for ETCS/ATP.
  \item \textbf{$N = 2$} consecutive failures (fixed) \arr\ \texttt{DEGRADED}.
  \item Recovery: new service session declared with same \svcref\ + signed token.
\end{itemize}

\subsection{NAT Addressing Model}
\label{sec:nat}

\textbf{ASAS address rule (resolved)}: post-NAT source IP of first service session
SIP signalling (SD\_PUB\_IP, RT-0097 §6.2.2.4). No SIP header extraction required.
Port configurable in FROP.

% ── 3. TERMINOLOGY REPLACEMENT ──────────────────────────────────────────────
\section{Terminology Replacement Rule in ETSI TSs}

\begin{longtable}{p{5.5cm} p{4cm} p{4.5cm}}
  \toprule Context in TSs & Current term & Proposed term \\ \midrule \endhead
  Internal MCx bearer & session & \textbf{service session} \\
  Internal MCx establishment & session establishment & \textbf{service session} establishment \\
  Internal MCx release & session release & \textbf{service session} release \\
  Inter-GW multi-path (service session level) & Multipath session & \textbf{Multi Service Path (MSP)} \\
  Transport-level Multipath (RT-00104) & Multipath & \textbf{Multipath} (unchanged) \\
  \OBAPP/\TSAPP\ identifier & \texttt{sessionId} & \textbf{unchanged} \\
  Application request & \texttt{POST /sessions}, \texttt{openSession*} & \textbf{unchanged} \\
  Incoming notification & \texttt{incomingSessionNotif*} & \textbf{unchanged} \\
  Application release & \texttt{closeSession*}, \texttt{DELETE /sessions} & \textbf{unchanged} \\
  Application state & \texttt{open session} & \textbf{unchanged} \\
  \bottomrule
\end{longtable}

% ── 4. PROPOSED ETSI MODIFICATIONS ──────────────────────────────────────────
\section{Proposed ETSI Modifications}

\subsection{ETSI TS 103\,764 (RT-0095) -- Normative Annex X}

\paragraph{X.1 -- Scope}
ASAS inter-GW API, agnostic to service session type. Not visible to applications.
Does not modify any UIC specification or service protocol. Absorbs
\texttt{FUTURE\_CRITICAL\_SERVICE} without modification of this annex.

\paragraph{X.2 -- Transport and Authentication}

\textbf{Transport:} HTTP/2 + JSON. Port: FROP. HTTP transport MAY be protected by
TLS~1.3 (recommended); there is no normative requirement for confidentiality of ASAS
exchanges between Gateways.

\textbf{Server certificate (if TLS):} Self-signed or any locally-deployed CA --
no inter-operator PKI required.

\textbf{ASAS IP address:} Post-NAT source IP of first service session SIP signalling
(SD\_PUB\_IP, RT-0097 §6.2.2.4). No SIP header extraction required.

\textbf{Normative authentication:} First MCx service session anchor + Ed25519 keys
exchanged over that trusted channel. All subsequent ASAS exchanges authenticated by
Ed25519 signature. No inter-GW PKI required.

\textbf{Token serialisation:} UTF-8 fields concatenated with \texttt{0x00}
separators (see §\ref{sec:trust}).

\textbf{Timers:} recovery timer INTEGER(1..180)\,s (FROP);
T\_KA INTEGER(1..10)\,s (application profile); $N = 2$ (fixed).

\textbf{mTLS option:} If bilateral OOB FSNNI agreement (FROP).

\paragraph{X.3 -- Initialisation (\texttt{POST /asas/appsessions})}

\begin{lstlisting}[language={}]
// Request
{ "appSessionId_call": "uuid-v4", "svcSessionRef_1": "uuid-v4",
  "svcSessionId_1": "native-id", "callingPublicKey": "base64(Ed25519)",
  "keepaliveReq": 2,
  "mspCapabilities": { "supportedPolicies": ["DUPLICATION","RECOVERY","LOADBALANCE"],
                       "maxPaths": 4 } }

// Response 200 OK
{ "appSessionId_term": "uuid-v4", "svcSessionRef_1": "uuid-v4",
  "terminatingPublicKey": "base64(Ed25519)", "keepaliveAccepted": 2,
  "mspPolicy": { "policy": "DUPLICATION|RECOVERY|LOADBALANCE|NONE", "activePaths": 2 },
  "advertisedAliases": [{ "svcSessionRef": "uuid-v4",
    "serviceType": "MCX|FUTURE_CRITICAL_SERVICE",
    "serviceDomain": "string", "serviceId": "string",
    "reachability": "CURRENT|ALTERNATIVE" }] }
\end{lstlisting}

\paragraph{X.4 -- Service Session Declaration / Recovery}

\begin{lstlisting}[language={}]
// Request
{ "appSessionId_call": "uuid-v4", "appSessionId_term": "uuid-v4",
  "svcSessionRef": "uuid-v4", "serviceType": "MCX|FUTURE_CRITICAL_SERVICE",
  "svcSessionId": "native-id", "svcSessRequestTime": "ISO8601",
  "svcSessAcceptTime": "ISO8601",
  "sessionToken": "base64(Sign(sk_call, quintet_0x00_separated))" }
// Recovery: add "replSvcSessionRef": "uuid-v4"

// Response 200 OK
{ "svcSessionRef": "uuid-v4",
  "ackToken": "base64(Sign(sk_term, same input))" }
\end{lstlisting}

\paragraph{X.5 -- Keepalive}
Sent only if applicative silence $\geq$ T\_KA on listed \svcref.

\begin{lstlisting}[language={}]
{ "appSessionId_call": "uuid-v4", "appSessionId_term": "uuid-v4",
  "svcSessionRefs": ["uuid-v4"], "timestamp": "ISO8601" }
// Response: { "results": [{ "svcSessionRef": "uuid-v4", "status": "ACTIVE|DEGRADED" }] }
\end{lstlisting}

\paragraph{X.6 -- Alias List Update (\texttt{PATCH .../aliases})}
Sent \textbf{immediately} by terminating GW on any new alias or service identifier
registration during an active session. No FROP configuration required.

\begin{lstlisting}[language={}]
{ "appSessionId_call": "uuid-v4", "appSessionId_term": "uuid-v4",
  "updatedAliases": [{ "svcSessionRef": "uuid-v4",
    "serviceType": "MCX|FUTURE_CRITICAL_SERVICE",
    "serviceDomain": "string", "serviceId": "string",
    "reachability": "CURRENT|ALTERNATIVE", "action": "ADD|REMOVE" }] }
\end{lstlisting}

\paragraph{X.7 -- Release (\texttt{DELETE .../appsessions/\{id\}})}

\begin{lstlisting}[language={}]
{ "appSessionId_call": "uuid-v4", "appSessionId_term": "uuid-v4",
  "reason": "NORMAL|ERROR|TIMEOUT" }
\end{lstlisting}

\paragraph{X.8 -- State Machines}

\textbf{Session:}
\begin{lstlisting}[language={}]
[CREATING]    -- 200 OK           --> [ESTABLISHED]
[CREATING]    -- 4xx/5xx          --> [FAILED]
[ESTABLISHED] -- all svc lost     --> [DEGRADED]
[ESTABLISHED] -- DELETE /sessions --> [RELEASING]
[DEGRADED]    -- >=1 ACTIVE       --> [ESTABLISHED]
[DEGRADED]    -- timer expired    --> [RELEASING]
[RELEASING]   -- done             --> [RELEASED]
\end{lstlisting}

\textbf{Service session:}
\begin{lstlisting}[language={}]
[CREATING]        -- bearer+ASAS OK    --> [ACTIVE]
[ACTIVE]          -- 2 KA no response --> [DEGRADED]
[DEGRADED]        -- recovery          --> [ACTIVE]
[ACTIVE|DEGRADED] -- RELEASING         --> [RELEASED]
\end{lstlisting}

\paragraph{X.9 -- Error Mapping to \OBAPP\ Causes}

\begin{longtable}{p{5.5cm} p{8cm}}
  \toprule ASAS cause & \texttt{openSessionFinalAnswerNotif} cause \\ \midrule \endhead
  ASAS timeout & \texttt{MCX\_ENDPOINT\_NOT\_REACHABLE} \\
  ASAS 403 & \texttt{REMOTE\_ENDPOINT\_DECLINED} \\
  ASAS 404 & \texttt{TERMINATING\_APPLICATION\_ENDPOINT\_NOT\_REACHABLE} \\
  ASAS 401 & \texttt{TERMINATING\_APPLICATION\_ENDPOINT\_NOT\_ALLOWED} \\
  \bottomrule
\end{longtable}

% ── 5. IMPACT SUMMARY ───────────────────────────────────────────────────────
\section{Impact Summary}

\begin{longtable}{p{5cm} p{6.5cm} p{2cm}}
  \toprule Document & Nature of modification & UIC specs \\ \midrule \endhead
  \textbf{ETSI TS 103\,764} & Session + service session + AIR + MSP + ASAS +
    service identifier (§3.1) + ASM + \textbf{Annex X} & \\
  \textbf{ETSI TS 103\,765-3} & Systematic replacement + Step 1
    (AIR \arr\ service id + SD\_PUB\_IP) + MSP policy & unchanged \\
  \textbf{ETSI TS 103\,765-4} & Systematic replacement + Step 0 +
    immediate alias list mgmt + MSP & unchanged \\
  \textbf{ETSI TS 103\,765-5} & Systematic replacement + MSP binding +
    Multipath/MSP distinction & unchanged \\
  \textbf{UIC FFFIS-7950} & \textbf{No modification} & \cmark \\
  \textbf{UIC FIS-7970} & \textbf{No modification} & \cmark \\
  \textbf{UIC FRS FU-7120} & \textbf{No modification} & \cmark \\
  \textbf{UIC AT-7800 SRS} & \textbf{No modification} & \cmark \\
  \textbf{UIC TOBA FRS-7510} & \textbf{No modification} & \cmark \\
  \textbf{3GPP TS 24.582/24.379} & \textbf{No modification} & \\
  \bottomrule
\end{longtable}

% ── 6. OPEN POINTS ──────────────────────────────────────────────────────────
\section{Open Points}

All open points are resolved. The document is ready for ETSI contribution.

\begin{enumerate}[leftmargin=2em]
  \item[\st{1.}] ASAS IP / NAT -- \textit{Resolved} (§\ref{sec:nat}).
  \item[\st{2.}] FSNNI multi-hop trust -- \textit{Resolved} (§\ref{sec:trust}).
  \item[\st{3.}] sessionToken serialisation -- \textit{Resolved} (§\ref{sec:trust} + §X.2).
  \item[\st{4.}] AIR rules -- \textit{Resolved}: \texttt{*.ertms}
    \arr\ \texttt{@current\_domain}; otherwise pass-through.
  \item[\st{5.}] Recovery timer -- \textit{Resolved}: INTEGER(1..180)\,s;
    $\leq 5$\,s for ETCS/ATP.
  \item[\st{6.}] T\_KA / N -- \textit{Resolved}: INTEGER(1..10)\,s; $N=2$;
    T\_KA $\leq$ timer/2.
  \item[\st{7.}] Alias list triggering -- \textit{Resolved}: immediate, no FROP.
  \item[\st{8.}] Extensibility -- \textit{Resolved}: \texttt{serviceType} open field.
\end{enumerate}

\end{document}